\documentclass[11pt]{article}

\usepackage[margin=1in]{geometry}
\usepackage{amsmath, amssymb}
\usepackage{algorithm}
\usepackage{algpseudocode}
\usepackage{booktabs}
\usepackage{enumitem}
\usepackage{hyperref}
\usepackage{tikz}
\usetikzlibrary{positioning, arrows.meta}

\hypersetup{colorlinks=true, linkcolor=blue!50!black, urlcolor=blue!50!black, citecolor=blue!50!black}

\newcommand{\R}{\mathbb{R}}
\newcommand{\code}[1]{\texttt{#1}}

\title{SBG: A CityJSON Pipeline for Town-Scale Building Geometry\\
\large Backend Logic, Phases 0--5}
\author{}
\date{}

\begin{document}
\maketitle

\begin{abstract}
This document describes the logic of SBG (Singapore Building Geometry), a pipeline that
turns raw open building footprint and elevation data into a spec-valid CityJSON dataset,
and ultimately into a watertight solid mesh suitable for meshing into a computational fluid
dynamics (CFD) domain for town-scale radionuclide plume dispersion modelling. It covers the
mathematics and algorithms behind each processing phase --- extrusion, height backfill,
domain cutout, terrain draping, editing, and solid export --- independent of the specific
libraries or source code used to implement them. Web interface and interactive tooling are
out of scope; this is a description of the backend's reasoning, not its code.
\end{abstract}

\tableofcontents

\section{Overview}

The pipeline consumes three heterogeneous open-data sources for Singapore --- a 2D building
footprint layer with inconsistent height attribution, a government building register with
authoritative but sparse measured heights, and a coarse (20\,m interval) contour layer --- and
produces a single coherent 3D dataset. All internal geometry uses EPSG:3414 (SVY21), the
official Singapore projected coordinate system, in metres; all inputs are reprojected once at
ingest and degrees are never carried past that boundary.

Six phases build on one another:

\begin{enumerate}[label=\textbf{Phase \arabic*.}, leftmargin=1.6cm]
  \item \textbf{Extrusion.} Turn each 2D footprint and a resolved scalar height into a
    Level-of-Detail~1 (LoD1) prismatic solid.
  \item \textbf{Height backfill.} Replace estimated heights with authoritative measured
    heights wherever a government register match exists, and keep the 3D geometry consistent
    with that correction.
  \item \textbf{Cutout.} Extract the subset of buildings fully inside an arbitrary
    user-specified domain boundary.
  \item \textbf{Terrain overlay.} Drape the domain's buildings and a locally-triangulated
    terrain surface onto each other so they share an exact, gap-free seam.
  \item \textbf{Editing.} Add or remove individual buildings from an existing dataset without
    reprocessing the whole file.
  \item \textbf{Solid export.} Fuse the domain's buildings and terrain into a single watertight
    solid and simplify it to a resolution matching the downstream CFD mesh.
\end{enumerate}

A recurring theme, revisited throughout, is that \emph{correctness at the seams} --- where one
building's geometry meets another's, or where a building meets the terrain --- is harder than
it looks, and several early, more ``obvious'' designs were abandoned once real data broke
their implicit assumptions. Where that happened, this document says so.

\section{Phase 1: Footprint Extrusion}
\label{sec:extrusion}

\subsection{Height resolution}

Every source footprint $b$ carries at most an OpenStreetMap height tag $h_{\text{osm}}(b)$ and
a storey count $\ell(b)$; both are frequently absent (in the raw data, 73\% of footprints have
no usable height tag at all). A building's resolved height $h(b)$ is a coalesce over three
tiers, from most to least trusted:

\begin{equation}
  h(b) = \begin{cases}
    h_{\text{osm}}(b) & \text{if } h_{\text{osm}}(b) > 0 \\[2pt]
    \ell(b) \cdot H_{\text{storey}} & \text{if } h_{\text{osm}}(b) = 0 \text{ and } \ell(b) > 0 \\[2pt]
    H_{\text{default}} & \text{otherwise}
  \end{cases}
  \label{eq:height-coalesce}
\end{equation}

with $H_{\text{storey}} = H_{\text{default}} = 3.2$\,m. Each building is tagged with which
tier resolved it (\code{osm}, \code{levels}, or \code{estimated\_default}), purely for later
traceability --- a diagnostic pass confirmed the zero-height buildings are scattered
island-wide with no geographic clustering, ruling out a systematic data gap in one area, and
the resulting tier tags let any later consumer filter out the least-trusted geometry if
needed.

\subsection{Extrusion into a solid}

A footprint is one or more polygon parts (a building split across disjoint parts becomes a
\emph{CompositeSolid} instead of a single \emph{Solid}); each part is an exterior ring
$r_0$ plus zero or more interior hole rings $r_1, \dots, r_k$. Given a resolved height $h(b)$
and a flat base elevation $z_0$ (taken as $0$ at this phase; Phase 4 generalises this to a
non-flat, terrain-following base), extrusion to a closed box is purely geometric:

\begin{algorithm}[H]
\caption{Footprint $\to$ LoD1 solid}
\label{alg:extrude}
\begin{algorithmic}[1]
\Function{ExtrudeSolid}{$\{r_0,\dots,r_k\}$, $z_0$, $z_1$}
  \State $F \gets \emptyset$ \Comment{set of faces, each a list of vertex indices}
  \For{each ring $r_i$}
    \State $B_i \gets$ vertices of $r_i$ at $z_0$
    \State $T_i \gets$ vertices of $r_i$ at $z_1$
  \EndFor
  \State $F \gets F \cup \{\,\mathrm{reverse}(B_0), \dots, \mathrm{reverse}(B_k)\,\}$
    \Comment{bottom cap, wound opposite the top so its normal points down}
  \State $F \gets F \cup \{\,T_0, \dots, T_k\,\}$ \Comment{top cap}
  \For{each ring $r_i$}
    \For{each edge $(p_j, p_{j+1})$ of $r_i$, cyclically}
      \State $F \gets F \cup \{\,[\,(p_j,z_0),\,(p_{j+1},z_0),\,(p_{j+1},z_1),\,(p_j,z_1)\,]\,\}$
        \Comment{one wall quad per edge}
    \EndFor
  \EndFor
  \State \Return closed solid with boundary $F$
\EndFunction
\end{algorithmic}
\end{algorithm}

The bottom cap is wound in reverse relative to the top cap purely so that, under a consistent
right-hand winding-to-normal convention, the bottom face's normal points downward and the top
face's upward --- both faces then correctly point outward from the solid.

\subsection{Vertex pooling and quantization}

Adjacent buildings, and a building's own top/bottom/wall faces, frequently reference
geometrically identical points (a wall's bottom-left corner is the same point as the bottom
cap's corresponding corner). Rather than storing each occurrence independently, every vertex
insertion is routed through a pool that deduplicates by rounded coordinate:

\begin{equation}
  \mathrm{key}(x,y,z) = \big(\,\mathrm{round}(x,p),\ \mathrm{round}(y,p),\ \mathrm{round}(z,p)\,\big), \qquad p = 3 \text{ decimal places}
\end{equation}

with the pool holding a hash map from key to index; a repeated key returns the existing index,
a new key is appended and assigned the next index. This deduplication happens \emph{globally}
across the whole file being built, not per building --- a detail that matters again in
Section~\ref{sec:backfill-geometry}. Once the full vertex list is known, it is quantized to
integer coordinates for the CityJSON output:

\begin{equation}
  q(x) = \mathrm{round}\!\left(\frac{x - x_{\min}}{s}\right), \qquad s = 10^{-3}
  \label{eq:quantize}
\end{equation}

i.e.\ millimetre-precision integers relative to the dataset's own minimum corner, which
CityJSON's \texttt{transform} object records so any reader can dequantize
($x = q(x)\cdot s + x_{\min}$) exactly.

\section{Phase 1.5: Authoritative Height Backfill}
\label{sec:backfill}

\subsection{Nearest-centroid spatial join}

A separate, independently surveyed building register provides government-measured heights for
most, but not all, buildings, indexed only by an approximate point location, not by matching
footprint geometry. Matching a footprint $b$ to a register record is a nearest-neighbour
query: compute each footprint's centroid $c(b)$, build a $k$-d tree over every register
record's location, and for each $b$ find

\begin{equation}
  r^*(b) = \operatorname*{argmin}_{r \in \mathcal{R}} \lVert c(b) - \mathrm{loc}(r) \rVert
\end{equation}

accepting the match only if $\lVert c(b) - \mathrm{loc}(r^*(b)) \rVert \le \tau$ for a fixed
tolerance $\tau = 40$\,m; otherwise $b$ keeps whatever height Equation~\ref{eq:height-coalesce}
already gave it. In practice this matches 93\% of buildings, dropping the
\code{estimated\_default} tier from 73\% of the dataset to 6\%. Register records are
deduplicated by their own stable identifier before the tree is built, since the underlying
tile scheme yields the same physical building many times over at different levels of detail.

\subsection{Keeping geometry consistent with a height correction}
\label{sec:backfill-geometry}

Overwriting $h(b)$ alone is insufficient: the solid built in Section~\ref{sec:extrusion} has
its top face's vertices already placed at the \emph{old} height. An early version of this
phase did exactly that --- updated the attribute, left the geometry untouched --- and the
resulting mismatch was only caught by a user visually inspecting the rendered model and
noticing a well-known 58\,m building rendered flat at the 3.2\,m default. A systematic check
afterwards found the geometry silently disagreed with the (correct) attribute for effectively
every corrected building.

The naive fix --- re-run Algorithm~\ref{alg:extrude} from scratch for every corrected building
--- is correct but wasteful: it allocates a full new set of vertices even though, in the
common case, a building's own top-face and wall-top vertices are not shared with anything
else and can simply be \emph{moved}. But the deduplicating pool from
Section~\ref{sec:extrusion} operates globally, not per building, so two structurally attached
buildings (e.g.\ a row of shophouses sharing a party wall) that happen to have landed on the
same default height can genuinely share a vertex at their common top corner. Moving such a
vertex in place would silently drag a neighbour's roof along with it.

The correct approach computes, once, a global reference count for every vertex index (how many
times it is used across every solid in the file), then treats each corrected building's own
top vertices individually:

\begin{algorithm}[H]
\caption{Relocate a building's roof to a new height, without corrupting shared geometry}
\label{alg:reheight}
\begin{algorithmic}[1]
\Function{RegenerateHeight}{$b$, $h_{\text{new}}$, $\mathrm{refcount}[\cdot]$}
  \State $L \gets$ multiset of vertex indices referenced anywhere in $b$'s own solid
  \State $V_{\text{top}} \gets \{\, v \in L : z(v) \neq z_0 \,\}$ \Comment{top-face \& wall-top vertices}
  \State $\mathrm{map} \gets \{\, v \mapsto v : v \in L \,\}$ \Comment{identity by default}
  \For{$v \in V_{\text{top}}$}
    \If{$\mathrm{refcount}[v] > \mathrm{count}(L, v)$} \Comment{used outside $b$ too}
      \State $v' \gets$ new pool entry at $(x(v), y(v), h_{\text{new}})$
      \State $\mathrm{map}[v] \gets v'$
      \State $\mathrm{refcount}[v] \gets \mathrm{refcount}[v] - \mathrm{count}(L, v)$
    \Else \Comment{exclusively $b$'s own vertex}
      \State mutate the pool entry for $v$ in place: $z(v) \gets h_{\text{new}}$
    \EndIf
  \EndFor
  \State rewrite $b$'s solid boundary by applying $\mathrm{map}$ to every referenced index
\EndFunction
\end{algorithmic}
\end{algorithm}

The comparison $\mathrm{refcount}[v] > \mathrm{count}(L, v)$ --- global usage count strictly
exceeding \emph{this building's own} usage count of the same vertex --- is what distinguishes
``shared with someone else'' from ``referenced several times within my own solid,'' which the
naive check $\mathrm{refcount}[v] > 1$ cannot: a building's own top ring and its adjoining wall
quads legitimately reference the same corner vertex multiple times purely as an artefact of
how Algorithm~\ref{alg:extrude} lists faces. Run against a real 118{,}782-building dataset,
this reduced the redundant vertex growth from what a full re-extrusion would have cost
(roughly a 100\% increase) to 3.6\%, while a direct audit confirmed zero unrelated buildings
had their geometry altered.

One further correctness detail: the decision of \emph{whether} a building needs regeneration
at all must be made by comparing the candidate new height against the building's
\emph{current geometry}, not its current attribute --- if this phase is ever re-run against an
already-corrected file (for instance because the register itself was updated), comparing
against the attribute would find ``no change'' for every already-corrected building and skip
exactly the geometry that most needs re-checking.

\subsection{A more ambitious, later-abandoned approach: full mesh replacement}

Height correction fixes buildings whose \emph{footprint} is right but whose \emph{height} was
wrong. It cannot fix buildings whose real form is fundamentally not a uniform-height prism ---
a stadium's swept roof, a building's outline formed by unioning several offset stacked blocks
into one OSM polygon. For a short list of such landmark buildings, a separate effort attempted
to replace their LoD1 box with real surveyed mesh geometry extracted from a government 3D
tile service, transforming each mesh from its native East-North-Up frame (centred at a
per-tile reference point) through Earth-centred Earth-fixed coordinates into the pipeline's
own projected system.

This was built and technically worked --- coordinate accuracy was independently verified by
recovering a known building's exact height on the first attempt --- but was reverted after
visual review placed the replacement meshes \emph{next to their real neighbours} rather than
in isolation: several replacements were offset from their OneMap reference point by 70--300\,m,
and one showed two adjacent landmark meshes overlapping each other. The root cause was
determined to be a \emph{scope} mismatch, not a coordinate error: the OSM footprint being
replaced sometimes spans a whole podium-and-towers complex while the corresponding surveyed
mesh covers only the roof structure, or vice versa. Because an isolated, no-context render had
been used as the only validation step, this mismatch was invisible until the final review ---
a lesson recorded and carried forward into every later phase, all of which validate results in
their real surrounding context rather than in isolation wherever practical.

\section{Phase 2: Domain Cutout}
\label{sec:cutout}

Given an arbitrary polygon domain $D$ specified by a user (in practice, drawn as a boundary on
a map), the cutout operation keeps only buildings fully enclosed by $D$:

\begin{equation}
  \mathrm{keep}(b) \iff \mathrm{footprint}(b) \subseteq D
\end{equation}

with no partial clipping: a building straddling the boundary is dropped in its entirety rather
than being cut in half, since a half-building is not a meaningful CFD obstacle. This requires
an exact geometric containment test, not merely a bounding-box overlap test, since a domain
boundary is rarely axis-aligned.

The reverse of the vertex-pooling step in Section~\ref{sec:extrusion} is needed here: once the
kept building set $K$ is known, the retained CityJSON must reference a \emph{compacted}
vertex array (indices $0, \dots, |V_K|-1$) rather than sparse indices into the original file's
full array.

\begin{algorithm}[H]
\caption{Compact a subset of a CityJSON's geometry to a dense vertex range}
\begin{algorithmic}[1]
\Function{CompactVertices}{$K$}
  \State $U \gets \emptyset$ \Comment{referenced index set}
  \For{$b \in K$}
    \State $U \gets U \cup \{\text{every vertex index appearing in } b\text{'s boundary}\}$
  \EndFor
  \State $U_{\text{sorted}} \gets \mathrm{sort}(U)$
  \State $\mathrm{map} \gets \{\, U_{\text{sorted}}[i] \mapsto i : i = 0,\dots,|U|-1 \,\}$
  \State new vertex array $\gets [\, V[u] : u \in U_{\text{sorted}} \,]$
  \State rewrite every $b \in K$'s boundary by applying $\mathrm{map}$
\EndFunction
\end{algorithmic}
\end{algorithm}

This same compaction routine is reused, unmodified, by the editing operations in
Section~\ref{sec:editing} and by the height-regeneration step in
Section~\ref{sec:backfill-geometry}'s remap --- it is a general operation on any nested
boundary structure, independent of what the geometry represents.

For interactive use at full-dataset scale, a naive implementation of $\mathrm{keep}(b)$ ---
reconstructing every building's footprint polygon from its solid geometry on every query ---
is far too slow (measured at tens of seconds against a 118{,}782-building file for a single
query). The mitigation is a one-time spatial index: an R-tree over every footprint, built once
when the dataset is loaded, so a containment query only reconstructs the geometry of buildings
whose bounding box could plausibly intersect $D$ in the first place, rather than every
building in the file.

\section{Phase 3: Terrain Overlay}
\label{sec:terrain}

\subsection{Elevation surface from sparse contours}

The source elevation data is a sparse point cloud sampled along 20\,m-interval contour lines,
not a regular grid. A digital terrain model (DTM) --- a regular grid of elevations
$Z[i,j]$ --- is built by scattered-data interpolation: for output grid coordinates $(x,y)$,

\begin{equation}
  Z(x,y) = \begin{cases}
    \mathrm{LinearInterp}(x,y) & \text{if } (x,y) \text{ is inside the input points' convex hull} \\
    \mathrm{NearestInterp}(x,y) & \text{otherwise}
  \end{cases}
\end{equation}

using a Delaunay triangulation of the input point cloud for the linear case (each output cell
falls inside some input triangle and is interpolated from that triangle's three corner
elevations) and nearest-neighbour fill outside the hull, where linear interpolation is
undefined. A single global distance-to-nearest-sample threshold was tried, to suppress the
long, geometrically implausible ``flat wedge'' triangles that a Delaunay triangulation
produces when bridging a real coverage gap (for instance, mainland to a distant offshore
island) --- but this was abandoned: a single threshold cannot distinguish a genuine data gap
from Singapore's ordinary sparse contour spacing on flat terrain, and in testing it discarded
90\% of the entire island's grid as ``untrustworthy.'' Properly solving this would require a
real coastline mask or a concave-hull (alpha-shape) approach; the coarse interpolation is used
as-is, which is an acceptable approximation given this pipeline's town-scale accuracy
requirement.

Given a built DTM, elevation at an arbitrary point is a standard bilinear lookup: map
$(x,y)$ to fractional grid coordinates $(c,r)$ via the grid's affine transform, then

\begin{equation}
  Z(c,r) = (1-t_c)(1-t_r)\,Z_{00} + t_c(1-t_r)\,Z_{10} + (1-t_c)t_r\,Z_{01} + t_c t_r\,Z_{11}
\end{equation}

where $Z_{00},\dots,Z_{11}$ are the four grid cells surrounding $(c,r)$ and $t_c, t_r$ are the
fractional offsets within that cell.

\subsection{The seam problem, and why a per-vertex offset is not the fix}

Draping buildings onto a terrain surface independently of the terrain's own triangulation ---
sampling $Z$ at each building's centroid and shifting the whole prism by that one offset ---
is simple but wrong at any location where elevation varies meaningfully across a single
footprint: a diagnostic pass found 8.9\% of buildings in a real test domain have over 1\,m of
elevation variation across their own footprint, concentrated at terrain escarpments. Buildings
in such areas visibly float above or sink into the terrain, with a real gap at the
building/ground boundary. Sampling elevation at every footprint vertex individually (instead
of once at the centroid) reduces, but does not eliminate, the gap: the terrain surface and the
building's raked base are still two independently triangulated surfaces that merely pass near
each other, not surfaces that share vertices.

\subsection{Constrained triangulation}

The seam is closed \emph{by construction}, not by numerical proximity, by inserting every
building footprint boundary as a hard constraint directly into the terrain's own triangulation,
rather than triangulating the terrain independently and draping buildings onto the result
afterwards. This uses a planar straight-line graph (PSLG) --- a set of vertices, required
edges between some of them, and hole regions that must not be triangulated:

\begin{itemize}
  \item A background grid of terrain sample points fills the domain, excluding points inside
    (or within a small collar of) any building footprint.
  \item The \emph{union} of all footprints' boundaries --- not each building's boundary
    independently --- is inserted as required edges. Using the union rather than per-building
    boundaries means two touching or overlapping buildings contribute one shared boundary
    segment rather than two coincident ones, which avoids producing degenerate sliver
    triangles at their shared wall.
  \item Each footprint's interior holes (courtyards) are \emph{not} punched out of the terrain
    --- the ground continues to exist underneath a courtyard, so only the annulus between a
    footprint's outer and inner ring is treated as excluded, not the courtyard interior itself.
  \item The overall domain boundary is inserted as required edges too: omitting it was found,
    empirically, to make the triangulator silently ignore most of the background grid instead
    of raising an error, because it could no longer determine which region was ``inside.''
\end{itemize}

A constrained Delaunay triangulation of this PSLG then produces a terrain mesh whose vertex set
already includes a vertex at every point where a footprint touches the ground.
Building base vertices are then generated by evaluating the \emph{same} elevation function at
the \emph{same} $(x,y)$ coordinates that the terrain triangulation used at that boundary, so
they quantize (Equation~\ref{eq:quantize}) to the identical vertex-pool key and collapse to a
single shared index automatically --- the seam closes as a direct consequence of both surfaces
sampling elevation identically at identical points, not as an approximation. This was verified
numerically on real data (not merely asserted): every building base vertex on a footprint's
exposed boundary was confirmed to share the exact pool index of the corresponding terrain
vertex, with zero exceptions once vertices on \emph{hidden} shared walls between touching
buildings (which are legitimately absent from the terrain, since no ground is exposed there)
were correctly excluded from the check.

One numerical pitfall is worth recording: the triangulation library used silently produced
zero triangles when handed vertices at Singapore's real absolute coordinate magnitude
(tens of thousands of metres from the coordinate origin), for reasons of internal
floating-point precision handling, and produced a correct result once the same input was
translated to sit near the origin first. The fix is a simple shift-then-triangulate-then-shift-back:

\begin{equation}
  V' = V - \min(V), \qquad \text{triangulate}(V', E, H), \qquad V'_{\text{out}} = V_{\text{out}} + \min(V)
\end{equation}

applied uniformly to every vertex and hole seed point passed to the triangulator.

\subsection{Draped extrusion}

Once a footprint's base elevation is a function of position rather than a single scalar,
Algorithm~\ref{alg:extrude} needs one modification: the flat base $z_0$ becomes a per-vertex
function $z_0(x,y) = Z(x,y)$ (raked to follow the ground), while the roof remains a single
flat plane,

\begin{equation}
  z_1 = \max_{(x,y) \in \text{footprint boundary}} Z(x,y) \; + \; h(b)
\end{equation}

taking the \emph{maximum} base elevation across the footprint, not an average, so that the
roof never dips below any point of a raked base on sloped terrain (which would invert a wall
face). Since the base is no longer planar, each wall quad is no longer guaranteed to be a
planar face and is triangulated into two triangles rather than left as a quad.

Two usage regimes exist for this overlay: a coarse, whole-island pass appropriate at
118{,}782-building scale (using the simpler centroid-drape from
Section~\ref{sec:terrain}'s first subsection, accepting its known limitations, since a full
constrained triangulation is not needed at whole-island scale) and the constrained,
seam-exact version described above, used only at the scale of a single user-specified CFD
domain (hundreds to low thousands of buildings), where the seam quality directly matters for
mesh generation.

\section{Phase 4: Editing}
\label{sec:editing}

Adding a building to an existing dataset reuses Algorithm~\ref{alg:extrude} directly, with the
new building's ground elevation resolved through the same three-tier priority used throughout
the pipeline: an explicit flat elevation if the caller supplies one, otherwise the terrain
elevation function if one is available for the dataset (draped per vertex, exactly as in
Section~\ref{sec:terrain}), otherwise a flat $z=0$ fallback. Because the target file is already
finalized (its vertex array already quantized against a fixed \code{transform}), new vertices
are appended directly in quantized form rather than re-running the deduplicating pool used at
initial construction --- correct because an edit adds a small number of vertices relative to
the whole file, so the cost of the rare duplicate is negligible next to the cost of rebuilding
a hash index over a multi-million-vertex file for every single edit.

Removing a building is the reverse: delete its CityObject entry. Whether to run the vertex
compaction from Section~\ref{sec:cutout} afterwards is a genuine trade-off, not automatic: a
single ad-hoc removal leaves a handful of now-unreferenced vertices in the pool, which is
harmless (every downstream tool tolerates unused vertices) and far cheaper than an $O(|V|)$
compaction pass; compaction is offered as an explicit, opt-in operation for the case where many
edits have accumulated and file size matters.

\section{Phase 5: Watertight Solid Export}
\label{sec:solid}

The final phase fuses a domain's buildings and terrain --- so far two conceptually separate
kinds of CityJSON object, a set of closed solids and one open terrain sheet --- into a single
watertight solid appropriate for CFD volume meshing, then simplifies it to a resolution that
matches the downstream mesh's own resolution.

\subsection{Why an exact boolean union does not work}

The most direct approach --- compute the exact constructive-solid-geometry (CSG) union of
every building solid with a solidified terrain slab --- was tried first and abandoned. Batching
the union of roughly 2600 buildings against the terrain in a single operation catastrophically
fragmented the result (dozens of disconnected pieces, with parts of the terrain itself missing
from the output). Performing the unions one at a time instead of batching recovered most of
the geometry but still left hundreds of leftover fragments, traced to buildings whose footprint
has near-duplicate or collinear vertices --- an OSM data-quality artefact that also causes
exactly the kind of degenerate input that exact CSG solvers are known to handle poorly at the
coincident seams Section~\ref{sec:terrain} deliberately creates. An exact boolean solver is
simply not robust against this much real-world data irregularity at this object count.

\subsection{Implicit resampling instead of an exact union}

The terrain, as noted, is an open sheet, not a volume; it is first \emph{solidified} by
extruding it downward by a fixed depth to produce a genuine slab, and each building's base is
deliberately \emph{plunged} downward through that slab (rather than merely touching its top
surface) to guarantee real volumetric overlap instead of a knife-edge surface coincidence at
exactly the seam vertices Section~\ref{sec:terrain} worked to make numerically identical.

Everything is then simply \emph{joined} into one mesh --- with no boolean operation computed at
all, overlaps, self-intersections and all --- and passed through a \emph{voxel remeshing} step:
the input is resampled onto a regular 3D grid at a fixed resolution (an inside/outside
classification per grid cell, in the spirit of an implicit-surface or signed-distance
representation), and a single new, clean, self-consistent surface is re-extracted from that
grid. This step does not care whether its input is self-intersecting, overlapping, or
fragmented into hundreds of disconnected pieces, because it never attempts to compute an exact
topological union of the input faces at all --- it only classifies space. On a real test
domain, this produced a single connected body from input that would have defeated an exact
boolean solver.

The voxel resolution is chosen to match the resolution the downstream CFD mesh itself will
use; resampling geometry to a coarser triangle size than the mesh's own cell size discards no
real fidelity, since the CFD mesher would discard finer detail anyway. A voxel size finer than
necessary is not a free quality upgrade, either --- face count scales with surface area
divided by the voxel size squared, uniformly across the whole surface including large flat
walls and roofs that gain no real detail from finer resampling, so an unnecessarily fine voxel
size blows up downstream face count for no benefit.

A small number of degenerate boundary slivers (order tens of fragments, a handful of faces
each, sitting exactly on the domain's outer edge) remain after remeshing, a resampling
precision artefact rather than lost geometry; these are identified by splitting the result into
its disconnected loose parts and discarding every part below a small face-count threshold,
keeping exactly one surviving main body.

\subsection{Simplification by quadric error metric edge collapse}

Voxel remeshing samples uniformly by surface area, so a flat wall, roof, or terrain patch ends
up just as densely tessellated as any other region --- genuinely wasteful, since no new shape
information exists there. The mesh is simplified afterwards using \emph{quadric error metric}
(QEM) edge collapse, a standard mesh-simplification algorithm: every vertex $v$ accumulates a
symmetric $4\times4$ matrix $Q_v$ summing, over every face plane $p$ incident to $v$ (with unit
normal $n_p$ and offset $d_p$, i.e.\ plane equation $n_p^\top x + d_p = 0$), the outer product

\begin{equation}
  Q_v = \sum_{p \ni v} \begin{bmatrix} n_p n_p^\top & n_p d_p \\ d_p n_p^\top & d_p^2 \end{bmatrix}
\end{equation}

so that for any point $\bar{x}$ (in homogeneous coordinates), $\bar{x}^\top Q_v \bar{x}$ is
exactly the sum of squared distances from $\bar{x}$ to every plane incident to $v$. Collapsing
an edge $(v_1, v_2)$ to a single new point $\bar{v}$ costs

\begin{equation}
  \mathrm{cost}(v_1, v_2 \to \bar{v}) = \bar{v}^\top (Q_{v_1} + Q_{v_2})\, \bar{v}
\end{equation}

and the algorithm greedily collapses the globally cheapest available edge, recomputing the
cost of newly-adjacent edges after each collapse, until a target face count is reached. This
is simplification, not the earlier deep-extrude-and-union step: it is applied once, after the
mesh is already a single clean watertight-or-near-watertight body, purely to reduce
unnecessary triangle density.

A visually more targeted-sounding alternative --- collapsing only near-coplanar faces within an
angle threshold, which in principle should specifically simplify flat regions without touching
curved or detailed ones --- was tried and abandoned: on real data at this face count it did not
terminate in a reasonable time, a known characteristic of that class of algorithm on large
meshes, whereas quadric edge collapse simplified the same mesh by an order of magnitude in
single-digit seconds. Quadric collapse simplifies more uniformly rather than specifically
sparing flat regions, but the resulting volume was independently confirmed to match the
pre-simplification volume to within a fraction of a percent despite a substantial face-count
reduction, which was judged an acceptable trade for this pipeline's LoD1-box source geometry.

Edge-collapse simplification of this kind has an inherent sequential character --- collapsing
one edge changes the cost of every edge adjacent to the collapsed region, so successive
collapse decisions are not independent of one another, which is the reason this class of
algorithm resists straightforward parallelization. A separate, embarrassingly-parallel
simplification strategy (vertex clustering: merging every vertex that falls in the same cell
of a coarse spatial grid) exists and was considered, but was ruled out for this pipeline
specifically because it does not preserve the mesh's manifold structure, which would work
directly against the watertightness goal of this whole phase.

\subsection{Verifying and, if necessary, repairing watertightness}

A mesh is watertight when it is simultaneously edge-manifold (every edge borders exactly two
faces), vertex-manifold (the faces around every vertex form a single connected fan), and free
of self-intersections. Despite the seam-matching construction in Section~\ref{sec:terrain} and
the implicit resampling in this section, the result is not guaranteed to satisfy this exactly:
watertightness is verified directly rather than assumed, and only if that check fails is a
general-purpose mesh repair (closing residual small holes, one class of correctable
non-manifold configuration) invoked --- an expensive operation whose cost scales with input
face count, skipped whenever it would be redundant. In practice, simplification is run
\emph{before} this check, both because it removes genuine near-degenerate artefacts that make
the check more likely to pass on its own, and because deferring the (potentially necessary)
repair step until after the mesh is already an order of magnitude smaller keeps that step's
memory use within what commodity hardware can handle --- attempting repair on the
full-resolution pre-simplification mesh directly was measured to be both far slower and to
approach the memory ceiling of a typical single workstation, confirming this ordering is a
reliability requirement of the pipeline, not merely a speed convenience.

\section{Summary}

Table~\ref{tab:phases} summarises each phase's core operation and its key correctness
guarantee.

\begin{table}[h]
\centering
\begin{tabular}{@{}p{2.6cm}p{5.6cm}p{6cm}@{}}
\toprule
\textbf{Phase} & \textbf{Core operation} & \textbf{Correctness guarantee} \\
\midrule
1. Extrusion & Footprint + scalar height $\to$ LoD1 solid & Consistent winding; deduplicated, quantized vertices \\
1.5 Backfill & Nearest-centroid height correction & Geometry kept consistent with the corrected attribute, without disturbing unrelated shared vertices \\
2. Cutout & Whole-building domain filter & Exact containment test; compacted vertex range \\
3. Terrain & Constrained triangulation + draped extrusion & Building/terrain seam closes by shared vertex index, not proximity \\
4. Editing & Append or remove one building & Elevation resolved consistently with the dataset's own terrain model \\
5. Solid export & Voxel resample + QEM simplify + verify/repair & Single watertight, orientable solid at target mesh resolution \\
\bottomrule
\end{tabular}
\caption{Phase summary.}
\label{tab:phases}
\end{table}

Across all six phases, the recurring failure mode was the same: an approach that looked correct
in isolation --- a per-vertex height offset, an isolated render of one replacement mesh, an
exact boolean union, a global-only reference count --- broke specifically at the point where
one piece of geometry meets another. Every fix in this pipeline that survived contact with real
data replaced a locally-plausible approximation with a construction that makes the relevant
quantity (a seam, a shared vertex, a watertight boundary) exactly true by how the geometry is
built, rather than approximately true by tuning a tolerance.

\end{document}
